\section{总体分析}
本题的关键在于把连续运动、几何约束和路径参数放进同一套语言中统一处理。Q1与Q2共用盘入螺线和把手递推，差别只是Q2进一步引入了板凳实体宽度；Q3则把螺距判定改写为角点到板凳轴线距离的临界搜索；Q4把调头部分抽象为固定交点下的两段相切圆弧；Q5在固定Q4路径后，通过速度缩放反推可接受的头速上限。

因此，全文的建模顺序应当是“先路径、后约束、再搜索”。先通过螺线弧长和递推关系确定运动状态，再通过几何判定筛选可行区域，最后对螺距或速度上界进行数值求解。这个顺序也与题目给出的五个问题自然对应。
